\documentclass{amsart}

% 数学相关宏包
\usepackage{amsmath,mathtools,amsthm,amssymb}
\usepackage{bbm}
\usepackage{mathrsfs}
\usepackage{centernot}

% 图形和绘图
\usepackage{graphicx,float}
\usepackage{tikz,tikz-cd}
\usepackage{pgfplots}
\usepackage[framemethod=TikZ]{mdframed}
\usetikzlibrary{decorations.pathmorphing,positioning,arrows,automata,calc,shapes,patterns}
\pgfplotsset{compat=1.18}
\usepgfplotslibrary{statistics}

% 表格和列表
\usepackage{booktabs,multirow,colortbl}
\usepackage{enumitem}
\setlist[itemize,enumerate,description]{noitemsep}
\setlistdepth{9}

% 算法
\usepackage{algorithm,algorithmic}

% 字体和颜色
\usepackage{xcolor,listings}
\usepackage{xeCJK}
\setCJKmainfont{Noto Serif CJK SC}

% 页面和链接
\usepackage{fancyhdr,caption}
\usepackage{hyperref}
\usepackage{cleveref}

% 工具宏包
\usepackage{etoolbox,letltxmacro,docmute}
\usepackage{gensymb}

% 自定义设置
\renewcommand{\arraystretch}{1.5}
\let\originalmathbb\mathbb
\renewcommand{\mathbb}[1]{%
    \ifstrequal{#1}{1}{\mathbbm{#1}}{\originalmathbb{#1}}%
}

% 颜色定义
\definecolor{lightblue}{RGB}{173,216,230}
\definecolor{lightgreen}{RGB}{144,238,144}
\definecolor{lightyellow}{RGB}{255,255,224}
\definecolor{lightgray}{RGB}{211,211,211}

		% 超链接设置
		\hypersetup{
		    colorlinks=true,
		    linkcolor=red,
		    filecolor=magenta,
		    urlcolor=cyan,
		    pdftitle={\@title},
		    pdfauthor={\@author},
		    pdfsubject={数学笔记},
		    pdfkeywords={数学},
		    bookmarksnumbered=true,
		    bookmarksopen=true,
		    bookmarksopenlevel=6,
		    pdfstartview=FitH,
		    pdfpagemode=UseOutlines,
		    pdfborder={0 0 0}
		}

% 定理环境
\theoremstyle{plain}
\newtheorem{theorem}{Theorem}[section]
\newtheorem{lemma}[theorem]{Lemma}
\newtheorem{corollary}[theorem]{Corollary}
\newtheorem{proposition}[theorem]{Proposition}

\theoremstyle{definition}
\newtheorem{definition}[theorem]{Definition}
\newtheorem{example}[theorem]{Example}
\newtheorem{exercise}[theorem]{Exercise}
\newtheorem{problem}[theorem]{Problem}

\theoremstyle{remark}
\newtheorem{remark}[theorem]{Remark}
\newtheorem{note}[theorem]{Note}
\newtheorem{observation}[theorem]{Observation}

% 解答环境
\newtheorem*{solution}{Solution}
\newtheorem*{proof*}{Proof}

% 图片路径
\graphicspath{{F:/资料/2025-summer/figures/}}

\author{乐绎华, Yue Yihua}
\address{中山大学, Sun Yat-sen University}
\email{yueyh@mail2.sysu.edu.cn}
\date{\today}
\title{Mathematical Notes}

\begin{document}

\maketitle
\tableofcontents

Great question! Here's the \textbf{general structure} for proving theorems in Lean:


\section{\textbf{1. Basic Theorem Structure}}

\begin{lstlisting}[language=lean]
theorem theorem_name (parameters) : statement := by
  proof_tactics
\end{lstlisting}
\textbf{Components}:

\begin{itemize}
	\item \textbf{\lstinline{theorem}}: Keyword to declare a theorem
	\item \textbf{\lstinline{theorem_name}}: Name of your theorem
	\item \textbf{\lstinline{parameters}}: Input variables and hypotheses
	\item \textbf{\lstinline{statement}}: What you want to prove
	\item \textbf{\lstinline{by}}: Starts tactic mode
	\item \textbf{\lstinline{proof_tactics}}: Step-by-step proof using tactics
\end{itemize}


\section{\textbf{2. Common Proof Patterns}}

\subsection{\textbf{Pattern 1: Universal Quantification (\lstinline{∀})}}

\begin{lstlisting}[language=lean]
theorem my_theorem : ∀ x y : ℕ, x + y = y + x := by
  intros x y           -- Introduce variables
  exact Nat.add_comm x y  -- Use existing lemma
\end{lstlisting}
\subsection{\textbf{Pattern 2: Implication (\lstinline{→})}}

\begin{lstlisting}[language=lean]
theorem my_theorem (P Q : Prop) : P → (P → Q) → Q := by
  intros hP hPQ        -- Introduce hypotheses
  exact hPQ hP         -- Apply hPQ to hP
\end{lstlisting}
\subsection{\textbf{Pattern 3: Equality}}

\begin{lstlisting}[language=lean]
theorem my_theorem (a b c : ℕ) : a + (b + c) = (a + b) + c := by
  exact Nat.add_assoc a b c  -- Use associativity
\end{lstlisting}
\subsection{\textbf{Pattern 4: Existence (\lstinline{∃})}}

\begin{lstlisting}[language=lean]
theorem my_theorem : ∃ x : ℕ, x > 0 := by
  use 1                -- Provide witness
  norm_num             -- Prove 1 > 0
\end{lstlisting}

\section{\textbf{3. Step-by-Step Proof Strategy}}

\subsection{\textbf{Step 1: Understand the Goal}}

\begin{lstlisting}[language=lean]
theorem my_theorem : some_statement := by
  -- Look at: ⊢ some_statement
  sorry
\end{lstlisting}
\subsection{\textbf{Step 2: Break Down the Goal}}

Use tactics to simplify:

\begin{lstlisting}[language=lean]
theorem my_theorem : ∀ x y : ℕ, x > 0 → y > 0 → x + y > 0 := by
  intros x y hx hy     -- Now goal is: ⊢ x + y > 0
  sorry
\end{lstlisting}
\subsection{\textbf{Step 3: Use Available Tools}}

\begin{lstlisting}[language=lean]
theorem my_theorem : ∀ x y : ℕ, x > 0 → y > 0 → x + y > 0 := by
  intros x y hx hy
  -- Use existing lemma or build proof
  exact Nat.add_pos hx hy
\end{lstlisting}

\section{\textbf{4. Common Proof Tactics (Toolbox)}}

\subsection{\textbf{Goal Manipulation}}

\begin{itemize}
	\item \textbf{\lstinline{intros}}: Introduce variables/hypotheses
	\item \textbf{\lstinline{apply}}: Apply a lemma/theorem
	\item \textbf{\lstinline{exact}}: Provide exact proof term
	\item \textbf{\lstinline{rw}}: Rewrite using equations
	\item \textbf{\lstinline{simp}}: Simplify using known rules
\end{itemize}

\subsection{\textbf{Case Analysis}}

\begin{itemize}
	\item \textbf{\lstinline{cases}}: Split on constructors
	\item \textbf{\lstinline{induction}}: Proof by induction
	\item \textbf{\lstinline{by_cases}}: Split on boolean conditions
\end{itemize}

\subsection{\textbf{Existence}}

\begin{itemize}
	\item \textbf{\lstinline{use}}: Provide witness for \lstinline{∃}
	\item \textbf{\lstinline{obtain}}: Extract witness from \lstinline{∃}
\end{itemize}

\subsection{\textbf{Logic}}

\begin{itemize}
	\item \textbf{\lstinline{constructor}}: Split \lstinline{∧} goals
	\item \textbf{\lstinline{left/right}}: Choose side of \lstinline{∨}
	\item \textbf{\lstinline{exfalso}}: Prove from contradiction
\end{itemize}


\section{\textbf{5. Complete Example}}

\begin{lstlisting}[language=lean]
-- Theorem: If x divides y and y divides z, then x divides z
theorem dvd_trans (a b c : ℕ) : a ∣ b → b ∣ c → a ∣ c := by
  intros hab hbc           -- Introduce hypotheses
  -- hab : a ∣ b means ∃ k, b = a * k
  -- hbc : b ∣ c means ∃ m, c = b * m
  obtain ⟨k, hk⟩ := hab   -- Extract witness: b = a * k
  obtain ⟨m, hm⟩ := hbc   -- Extract witness: c = b * m
  use k * m                -- Claim: c = a * (k * m)
  calc c = b * m           := hm
       _ = (a * k) * m     := by rw [hk]
       _ = a * (k * m)     := by rw [Nat.mul_assoc]
\end{lstlisting}

\section{\textbf{6. Proof Styles}}

\subsection{\textbf{Tactic Mode} (Step-by-step)}

\begin{lstlisting}[language=lean]
theorem example1 : a + b = b + a := by
  rw [Nat.add_comm]
\end{lstlisting}
\subsection{\textbf{Term Mode} (Direct)}

\begin{lstlisting}[language=lean]
theorem example2 : a + b = b + a := 
  Nat.add_comm a b
\end{lstlisting}
\subsection{\textbf{Calc Mode} (Calculational)}

\begin{lstlisting}[language=lean]
theorem example3 : (a + b) + c = a + (b + c) := 
  calc (a + b) + c = a + (b + c) := Nat.add_assoc a b c
\end{lstlisting}

\section{\textbf{7. General Workflow}}

\begin{enumerate}
	\item \textbf{State the theorem}
	\item \textbf{\lstinline{intros}} to bring in parameters
	\item \textbf{Analyze the goal} (\lstinline{⊢ ...})
	\item \textbf{Choose appropriate tactics}:
	\begin{itemize}
		\item Apply existing lemmas (\lstinline{apply}, \lstinline{exact})
		\item Rewrite equations (\lstinline{rw})
		\item Case analysis (\lstinline{cases}, \lstinline{induction})
		\item Provide witnesses (\lstinline{use})
	\end{itemize}
	\item \textbf{Check intermediate goals}
	\item \textbf{Repeat until \lstinline{no goals}}
\end{enumerate}


\section{\textbf{8. Pro Tips}}

\begin{itemize}
	\item \textbf{Start simple}: Use \lstinline{sorry} placeholders, then fill them in
	\item \textbf{Use \lstinline{#check}}: Verify types of lemmas
	\item \textbf{Read the goal}: Always understand what \lstinline{⊢} shows
	\item \textbf{Chain tactics}: Use semicolons \lstinline{;} for multiple goals
	\item \textbf{Break complex proofs}: Use \lstinline{have} for intermediate steps
\end{itemize}


\section{\textbf{Template to Follow}}

\begin{lstlisting}[language=lean]
theorem my_theorem (assumptions) : conclusion := by
  -- Step 1: Introduce variables
  intros variables_and_hypotheses
  
  -- Step 2: Apply relevant lemmas or analyze cases
  apply some_lemma  -- or cases, induction, etc.
  
  -- Step 3: Handle remaining goals
  exact proof_term  -- or more tactics
\end{lstlisting}
This structure works for most mathematical proofs in Lean!

\section{How does lean understand the \lstinline{intros}}

Great question! Lean understands what variables mean through \textbf{type information} and \textbf{the structure of the goal}. Here's how it works:


\subsection{\textbf{1. Lean Reads the Goal Structure}}

When Lean sees a goal like:

\begin{lstlisting}[language=lean]
⊢ ∀ {R : Type v} [inst : Mul R] [inst_1 : Add R] [RightDistribClass R] (a b c : R), 
  (a + b) * c = a * c + b * c
\end{lstlisting}
It \textbf{parses the structure} from left to right:

\begin{enumerate}
	\item \lstinline{∀ {R : Type v}} → "There's an implicit type variable \lstinline{R}"
	\item \lstinline{[inst : Mul R]} → "There's a typeclass instance for multiplication on \lstinline{R}"
	\item \lstinline{[inst_1 : Add R]} → "There's a typeclass instance for addition on \lstinline{R}"
	\item \lstinline{[RightDistribClass R]} → "There's a typeclass instance for right distributivity on \lstinline{R}"
	\item \lstinline{(a b c : R)} → "There are three explicit variables of type \lstinline{R}"
	\item \lstinline{(a + b) * c = a * c + b * c} → "The conclusion to prove"
\end{enumerate}


\subsection{\textbf{2. Type Information Guides Understanding}}

Each variable has \textbf{explicit type annotations}:

\begin{lstlisting}[language=lean]
{R : Type v}           -- R is a type (implicit)
[inst : Mul R]         -- inst is a Mul typeclass instance (implicit)
[inst_1 : Add R]       -- inst_1 is an Add typeclass instance (implicit)
[RightDistribClass R]  -- unnamed instance of RightDistribClass (implicit)
(a b c : R)            -- a, b, c are elements of type R (explicit)
\end{lstlisting}

\subsection{\textbf{3. How \lstinline{intros} Processes This}}

When you use \lstinline{intros}, Lean works \textbf{sequentially}:

\begin{lstlisting}[language=lean]
intros R inst inst_1 rightDistrib a b c
\end{lstlisting}
\textbf{Step by step}:

\begin{enumerate}
	\item \textbf{\lstinline{R}} → Lean knows this should be a \lstinline{Type v} (from \lstinline{{R : Type v}})
	\item \textbf{\lstinline{inst}} → Lean knows this should be \lstinline{Mul R} (from \lstinline{[inst : Mul R]})
	\item \textbf{\lstinline{inst_1}} → Lean knows this should be \lstinline{Add R} (from \lstinline{[inst_1 : Add R]})
	\item \textbf{\lstinline{rightDistrib}} → Lean knows this should be \lstinline{RightDistribClass R} (from \lstinline{[RightDistribClass R]})
	\item \textbf{\lstinline{a b c}} → Lean knows these should be elements of type \lstinline{R} (from \lstinline{(a b c : R)})
\end{enumerate}


\subsection{\textbf{4. Automatic Name Generation}}

If you don't provide names:

\begin{lstlisting}[language=lean]
intros  -- No names provided
\end{lstlisting}
Lean generates names based on:

\begin{itemize}
	\item \textbf{Original names} (if available): \lstinline{R}, \lstinline{inst}, \lstinline{inst_1}
	\item \textbf{Type-based names}: \lstinline{a✝}, \lstinline{a✝¹}, \lstinline{a✝²} for generic variables
	\item \textbf{Pattern recognition}: Common patterns get predictable names
\end{itemize}


\subsection{\textbf{5. Example: Step-by-Step Breakdown}}

\begin{lstlisting}[language=lean]
theorem my_theorem : ∀ {R : Type v} [Mul R] [Add R] [RightDistribClass R] (a b c : R), 
  (a + b) * c = a * c + b * c := by
  -- Goal: ⊢ ∀ {R : Type v} [Mul R] [Add R] [RightDistribClass R] (a b c : R), (a + b) * c = a * c + b * c
  
  intro R
  -- Goal: ⊢ [Mul R] → [Add R] → [RightDistribClass R] → ∀ (a b c : R), (a + b) * c = a * c + b * c
  -- Context: R : Type v
  
  intro inst_mul
  -- Goal: ⊢ [Add R] → [RightDistribClass R] → ∀ (a b c : R), (a + b) * c = a * c + b * c
  -- Context: R : Type v, inst_mul : Mul R
  
  intro inst_add
  -- Goal: ⊢ [RightDistribClass R] → ∀ (a b c : R), (a + b) * c = a * c + b * c
  -- Context: R : Type v, inst_mul : Mul R, inst_add : Add R
  
  intro inst_distrib
  -- Goal: ⊢ ∀ (a b c : R), (a + b) * c = a * c + b * c
  -- Context: R : Type v, inst_mul : Mul R, inst_add : Add R, inst_distrib : RightDistribClass R
  
  intros a b c
  -- Goal: ⊢ (a + b) * c = a * c + b * c
  -- Context: R : Type v, ..., a b c : R
\end{lstlisting}

\subsection{\textbf{6. Key Insight: Type-Driven}}

Lean's understanding is \textbf{entirely type-driven}:

\begin{itemize}
	\item Variables get their \textbf{meaning} from their \textbf{type}.
	\item The \textbf{order} in the goal determines the \textbf{order} of introduction.
	\item \textbf{Names} are just labels—the \textbf{types} are what matter for correctness.
\end{itemize}


\subsection{\textbf{7. What If You Get It Wrong?}}

If you try to assign incompatible names:

\begin{lstlisting}[language=lean]
intros my_number my_function  -- Wrong types
\end{lstlisting}
Lean will give a \textbf{type error} because it expects specific types in specific positions.


\subsection{\textbf{Summary}}

Lean understands variables through:

\begin{enumerate}
	\item \textbf{Type annotations} in the goal (\lstinline{R : Type v}, \lstinline{a : R}, etc.)
	\item \textbf{Sequential processing} of quantifiers and implications
	\item \textbf{Type checking} to ensure consistency
	\item \textbf{Context building} as each variable is introduced
\end{enumerate}

The key is that \textbf{types carry all the semantic information}—names are just convenient labels!


\end{document}